\documentclass[tikz,border=15pt]{standalone}
\usepackage{ctex}
\usepackage{tikz}
\usetikzlibrary{arrows.meta, positioning, calc}

\begin{document}
\begin{tikzpicture}[
    font=\normalsize,
    stepbox/.style={draw, thick, rounded corners=3pt, fill=blue!8, inner sep=5pt, minimum width=2.6cm, minimum height=1.3cm, align=center},
    arrow/.style={thick, -{Stealth[length=2.5mm]}, blue!70!black},
    pathchar/.style={font=\bfseries\large, red},
    coordlabel/.style={font=\small\ttfamily},
    pdlabel/.style={font=\small, green!50!black}
]

% 标题
\node[font=\Large\bfseries] at (0,6.2) {回溯方案：从 $G$ 回到 $S$};
\node[font=\small, below=2pt] at (0,5.7) {$p[x][y][d]$ 记录到达状态 $(x,y,d)$ 时，上一个状态的出方向};

% 第一行
\node[stepbox] (s1) at (-5.2,4.2) {
\coordlabel 当前 $(3,5,0)$\\
\pdlabel $pd = p[3][5][0] = R$\\[2pt]
\pathchar path += 'R'
};

\node[stepbox] (s2) at (-2.2,4.2) {
\coordlabel 当前 $(3,4,0)$\\
\pdlabel $pd = p[3][4][0] = R$\\[2pt]
\pathchar path += 'R'
};

\node[stepbox] (s3) at (0.8,4.2) {
\coordlabel 当前 $(3,3,0)$\\
\pdlabel $pd = p[3][3][0] = D$\\[2pt]
\pathchar path += 'D'
};

\node[stepbox] (s4) at (3.8,4.2) {
\coordlabel 当前 $(2,3,1)$\\
\pdlabel $pd = p[2][3][1] = D$\\[2pt]
\pathchar path += 'D'
};

% 第二行
\node[stepbox] (s5) at (-5.2,1.5) {
\coordlabel 当前 $(1,3,1)$\\
\pdlabel $pd = p[1][3][1] = U$\\[2pt]
\pathchar path += 'U'
};

\node[stepbox] (s6) at (-2.2,1.5) {
\coordlabel 当前 $(2,3,3)$\\
\pdlabel $pd = p[2][3][3] = U$\\[2pt]
\pathchar path += 'U'
};

\node[stepbox] (s7) at (0.8,1.5) {
\coordlabel 当前 $(3,3,3)$\\
\pdlabel $pd = p[3][3][3] = R$\\[2pt]
\pathchar path += 'R'
};

\node[stepbox] (s8) at (3.8,1.5) {
\coordlabel 当前 $(3,2,0)$\\
\pdlabel $pd = p[3][2][0] = D$\\[2pt]
\pathchar path += 'D'
};

% 箭头连接
\draw[arrow] (s1.east) -- (s2.west);
\draw[arrow] (s2.east) -- (s3.west);
\draw[arrow] (s3.east) -- (s4.west);
\draw[arrow] (s4.south) -- ++(0,-0.5) -| (s5.north);
\draw[arrow] (s5.east) -- (s6.west);
\draw[arrow] (s6.east) -- (s7.west);
\draw[arrow] (s7.east) -- (s8.west);

% 终点S
\node[draw, thick, rounded corners=3pt, fill=green!15, inner sep=6pt, font=\bfseries] (send) at (6.2,2.8) {到达 $S$!};
\draw[arrow] (s4.east) -- ++(0.6,0) |- (send.west);

% 回溯公式
\node[draw, thick, rounded corners=5pt, fill=yellow!15, inner sep=8pt, font=\small, text width=12cm, align=left] at (0,-1.0) {
\textbf{回溯公式：}\\[3pt]
\texttt{pd = p[x][y][d]} \quad $\Rightarrow$ \quad 上一个出方向\\
\texttt{px = x - dx[pd],\; py = y - dy[pd]} \quad $\Rightarrow$ \quad 上一步位置\\
\texttt{path.push\_back(dc[pd])} \quad $\Rightarrow$ \quad 记录路径字符\\[3pt]
循环直到 $(x,y) = (sx,sy)$，最后 \texttt{reverse(path)}
};

% 最终路径
\node[font=\large\bfseries, red] at (0,-2.8) {path = [R, R, D, D, U, U, R, D]};
\node[font=\large\bfseries, green!50!black] at (0,-3.3) {reverse(path) = D-R-U-U-D-D-R-R};

\end{tikzpicture}
\end{document}
